% =====================================================================
%  RESUMO EM LÍNGUA PORTUGUESA (obrigatório) -- Manual, seção 2.1.8
% ---------------------------------------------------------------------
%  Regras (ABNT NBR 6028:2021):
%   - UM ÚNICO parágrafo, sem recuo, espaço simples;
%   - entre 150 e 500 palavras;
%   - frases concisas e afirmativas; terceira pessoa do singular e voz
%     ativa ("Este trabalho analisa...", e não "Neste trabalho eu analisei...");
%   - deve apresentar: OBJETIVO, MÉTODO, RESULTADOS e CONCLUSÕES;
%   - sem citações, figuras, tabelas ou equações;
%   - palavras-chave logo abaixo, separadas por ponto e vírgula (;) e
%     terminadas com ponto final (.).
%  DICA: escreva o resumo por último, quando o trabalho estiver pronto.
% =====================================================================
\setlength{\absparsep}{18pt}
\begin{resumo}
\SingleSpacing
\noindent A evasão de praticantes é um problema recorrente em academias de
lutas e artes marciais, e muitos alunos abandonam o treino antes de
perceber os seus benefícios físicos e mentais. Este trabalho apresenta o
RoundUp, um aplicativo móvel gamificado que busca aumentar o engajamento e
a retenção dos alunos e apoiar professores e academias na gestão das
turmas. O projeto foi conduzido como um processo de ideação: treze ideias
de funcionalidades foram convertidas em requisitos, que orientaram a
construção do aplicativo e, ao mesmo tempo, foram avaliados por meio de um
questionário respondido por praticantes e ex-praticantes. Para ampliar a
cobertura dos perfis de usuário, incluindo o de professor, foram
empregadas personas sintéticas e um gerador de respostas desenvolvidos
com apoio de inteligência artificial: a partir de 12 personas, foram
produzidas 65 respostas que permitiram exercitar todo o questionário e o
procedimento de análise, mantidas separadas dos dados reais. O protótipo foi desenvolvido em React
Native com Firebase, com três perfis de usuário (aluno, professor e
academia), e implementa 29 dos 33 requisitos funcionais, entre eles
registro de presença com validação pelo professor, pontos de experiência,
níveis, sequência de treinos, medalhas, classificação, missões, registro
de graduação, relatórios de frequência e consentimento conforme a LGPD. As
regras de negócio que protegem a gamificação foram implementadas nas
regras de segurança do banco de dados e cobertas por testes automatizados.
Na avaliação real, os recursos que tornam o progresso visível, como a
barra de progresso até a graduação e o histórico de presença, tiveram as
melhores notas; lembretes e mensagens automáticas, as piores; e os
comentários apontaram o excesso de notificações e de dados a inserir como
motivos de abandono do aplicativo, o que levou a tornar as notificações
configuráveis e a classificação opcional. Conclui-se que a gamificação
centrada no progresso visível é a direção mais promissora para o
aplicativo, que parte da evasão tem causas externas que um aplicativo não
resolve e que os requisitos de gestão ainda precisam ser avaliados por
professores reais.

\vspace{\onelineskip}
\noindent\textbf{Palavras-chave}: gamificação; artes marciais; retenção; engenharia de requisitos; personas.
\end{resumo}
